\subsection{Mission Statement}

There are increasing security concerns in the Baltic region, primarily associated with ongoing conflicts. Ballistic missile launches pose a significant risk to civilian populations due to insufficient warning times. Therefore, national security authorities require a wide-coverage, high-reliability detection system capable of persistent monitoring and timely alerting.

We propose a reliable space-based satellite system to detect missile launches at the earliest possible stage and to track them sufficiently enough to support rapid, preliminary impact-area estimation. The system will distribute warning information to national authorities responsible for national crisis management and emergency response. Additionally, the system should include the capability to store historical data of previous events for algorithm improvement, system calibration, performance verification, operator training, and post-event technical analysis, in line with R-004.

\subsection{Mission Definition}

\subsubsection{Mission Need and Purpose}

The Early Warning Alarm System (\abbr{EWAS}) is defined as a space-based civil protection mission intended to provide national authorities responsible for national crisis management and emergency response with an early warning capability against ballistic missile launch threats. The mission is motivated by the limited warning time available during missile attacks and by the need to support rapid governmental and emergency-response decisions before a potential impact occurs.

The purpose of \abbr{EWAS} is to detect missile launch events at the earliest feasible stage, support preliminary assessment of the threat, estimate the potential impact area with an appropriate level of uncertainty, and deliver warning information to aforementioned authorities

\abbr{EWAS} is not intended to function as an offensive military system or as a weapon-control architecture. Its purpose is to provide timely warning and decision-support information for civil protection, national resilience, and emergency preparedness.

\subsubsection{Mission Objectives}

The \abbr{EWAS} mission objectives are to provide a reliable space-based capability that enables the detection, assessment and continuous improvement of missile early warning information.

The mission is intended to:

\begin{enumerate}
    \item Detect ballistic missile launch events affecting the Baltic region as early as possible.
    \item Track and process detected events sufficiently to support rapid preliminary estimation of the potential impact area.
    \item Generate clear warning information for authorized national crisis-management and emergency-response authorities.
    \item Maintain selected historical event data to improve future warning quality, support post-event analysis, refine classification methods, reduce false alarms, and calibrate preliminary impact-area estimation models.
    \item Provide actionable warning products to authorized national authorities for decisions concerning protection of the Baltic region, including potential threats originating outside that region.
\end{enumerate}

\subsubsection{Mission Scope and Boundaries}

The scope of the \abbr{EWAS} mission includes the space-based detection of missile launch events, preliminary event processing, initial threat assessment, warning generation, dissemination of warning products to authorized users, and storage of selected event data for technical and operational improvement.

Within this scope, \abbr{EWAS} should support:

\begin{itemize}
    \item Automated detection and preliminary event screening, followed by operator review of issued warnings;
    \item Preliminary classification of the detected event;
    \item Preliminary missile characterization based on observable event features and available aviation data;
    \item Preliminary trajectory and impact-area estimation;
    \item Secure delivery of warning information to national authorities responsible for national crisis management and emergency response;
    \item Post-event analysis and improvement of detection, classification, and estimation algorithms.
\end{itemize}

The mission boundaries are defined to preserve the civil protection character of the system. \abbr{EWAS} does not include missile interception, offensive targeting, weapon guidance, direct command of military response assets, final attribution of launch responsibility, or complete damage assessment. It also does not replace the national public alerting infrastructure; instead, it provides warning information to the authorities responsible for activating such systems.

\subsubsection{Definition of Mission Success}

Mission success for \abbr{EWAS} is defined by the system's ability to provide timely, reliable, and operationally useful missile warning information to national authorities responsible for civil protection and emergency response.

The mission may be considered successful if it can:

\begin{enumerate}
    \item Detect relevant missile launch events within the defined surveillance region.
    \item Generate a preliminary assessment of the event, including confidence level and estimated threat area.
    \item Maintain sufficient reliability and availability to support operation during crisis conditions.
    \item Provide warning products that are understandable and actionable for national emergency-management users without requiring them to operate or interpret the space system directly. Warning and data should be delivered within a time frame useful for crisis-management and civil-protection decisions.
    \item Store mission data securely and use historical event data only for the technical purposes permitted by R-004, including post-event review, calibration, training, verification, and refinement of classification and estimation methods.
\end{enumerate}

\subsubsection{High-Level System Boundary}
The current system baseline contains two system segments: the Space Segment and the Ground Segment. Governmental crisis-management and civil-protection authorities are external recipients of warning products rather than a separate User and Data Segment. Within the Ground Segment, the Early Warning System prepares and dispatches these products and the Data Archiving System retains the applicable records. This allocation follows the Ground Systems decomposition in Section~\ref{sec:ground-segment-boundary} and the existing work breakdown structure.

A launch may occur outside the Baltic region; the operational purpose is to assess whether the event could threaten the Baltic region and to provide appropriate warning information to its authorized authorities. This mission concept does not, by itself, demonstrate continuous global detection capability; supporting constellation-coverage evidence is to be supplied by the Space Segment team.
